\scanpage{55}
\mspage{23}
avec la structure de stratification cylindrique \inserted{ordonnée}
locale. \deleted{Par exemple}, Une autre façon de
décrire les $\Gamma(d)=\Sigma^{<}(d+1)$, c'est de
prendre la puiss. \deleted{\unsure{sym.}} \inserted{cart. relative} $d$.ième de $\Gamma_0$\deleted{\unsure{$(\varnothing)$}}\inserted{$=\Sigma^{<}(1)$}
sur $\Sigma\ (=\Gamma(\varnothing))$, dont on enlève les images
des puiss. cart. relatives \deleted{\unsure{inférieures}} \inserted{\deleted{$\Gamma'(d)$} des $d'$.ième ($d'<d$)} par les applic\unsure{ns}
$\tfrac12$ simpliciales entre celles ci. Cela donne à priori
les \deleted{\unsure{ens tot.}} suites ordonnées de $d$ germes de
strates de codim $1$ \inserted{en un $=$ point,} \unsure{(mutuel\textsuperscript{t}|\deleted{(mutuel\textsuperscript{t}}} distinctes), \deleted{soit}
$\Gamma'(d)$. On a alors \inserted{$\unsure{\Gamma^{\deleted{'}}(d)|\Gamma(d)}\subset\Gamma'(d)$},
$\Gamma(d)$ est une partie à la fois ouverte et
\unsure{fermée|\deleted{fermée}} de $\Gamma'(d)$, et \deleted{$\Gamma'(d)\unsure{\simeq}$}
\[
(35)\qquad \Gamma'(d)\simeq\unsure{\mathfrak S_d|\mathcal S_d}\times\Gamma(d)
\]
Pour avoir \deleted{la}\inserted{une} structure d'ordre \inserted{total} relative, de
\margin{\deleted{il s'agit de}}
\deleted{\unsure{$\Gamma'(1)$ sur}} $\Gamma_0=\Sigma^{<}(1)$ sur $\Sigma_{\deleted{\unsure{*|0}}}=\Sigma^{<}(0)$, il \deleted{suffit}\inserted{\unsure{uniquement}}
de se donner, \deleted{dans $\Gamma'_1(\unsure{0|2})=\unsure{\Sigma^{<}(0)|\Sigma^{<}(2)}$, $\Gamma'_{1+1}$} dans
$\Gamma_0\times_\Sigma\Gamma_0\setminus\text{diagonale}$, une partie \inserted{(i.e. \unsure{néc.}} ouverte et fermée)
$\Gamma_1$, telle que $\Gamma_1\cup{}^s\Gamma_1=\Gamma_0\times\Gamma_0\setminus\text{diag}$, $\Gamma_1\cap{}^s\Gamma_1=\varnothing$
\margin{$\Gamma_1$ est le graphe de la relation $x<y$}
i.e. une section de $\Gamma_0\times_\Sigma\Gamma_0\setminus\text{diag}$ \unsure{au|\deleted{au}} dessus
\margin{ou une section \deleted{de $\Gamma_0\times_\Sigma\Gamma_0\setminus\unsure{\text{diagonale}|\text{diag}}$}\deleted{de $\operatorname{Sym}^2(\Gamma_0\times_\Sigma\Gamma_0)$}\deleted{de}}
de $\operatorname{Sym}^2_\Sigma(\Gamma_0)^*$ (``choix d'un premier élément parmi deux
éléments distincts), avec la condition de transitivité
\deleted{$\unsure{\text{que }\Gamma_1}$}
\diagram{\begin{tikzcd}[column sep=large,row sep=small]
(36)\quad (\Gamma_1\times_\Sigma\Gamma_0)\cap(\Gamma_0\times_\Sigma\Gamma_1)
 \arrow[r,"{\operatorname{Pr}_{\unsure{13|23}}}"] \arrow[d,phantom,"\cap" description]
 & \Gamma_0\times\Gamma_0 \arrow[d,phantom,"\cup" description]
 & \text{se factorise par }\Gamma_1 \\
\unsure{\Sigma_0|\Gamma_0}\times\Gamma_0\times\Gamma_0 & \Gamma_1 &
\end{tikzcd}}
\margin{54}
